\begin{trapbox}

	\begin{description}[style=nextline, leftmargin=2.8em, labelsep=0.5em, itemsep=0.35em]
		\item[\textbf{\textcolor{CTrap}{易错点一（边长混淆）}}]
			容易将弦图大正方形的边长误认为 $a+b$。
		\item[\textbf{\textcolor{CTrap}{正确理解（边长归属）：}}]
			在证明勾股定理的标准弦图中，大正方形由斜边构成，边长为 $c$，小正方形边长为 $a-b$。
		\item[\textbf{\textcolor{CTrap}{错因分析（图形混淆）：}}]
			混淆了两种常见弦图拼法，错误地使用大正方形边长 $a+b$ 的关系。
	\end{description}

	\medskip
	\begin{description}[style=nextline, leftmargin=2.8em, labelsep=0.5em, itemsep=0.35em]
		\item[\textbf{\textcolor{CTrap}{易错点二（面积系数遗漏）}}]
			计算四个直角三角形总面积时，容易忽略 $\frac12$，写出 $4ab$。
		\item[\textbf{\textcolor{CTrap}{正确理解（系数注意）：}}]
			每个三角形面积 $\frac12 ab$，四个总面积为 $2ab$。
		\item[\textbf{\textcolor{CTrap}{错因分析（公式误用）：}}]
			三角形面积公式记忆错误，或误以为四个三角形面积即 $4ab$。
	\end{description}

	\medskip
	\begin{description}[style=nextline, leftmargin=2.8em, labelsep=0.5em, itemsep=0.35em]
		\item[\textbf{\textcolor{CTrap}{易错点三（条件忽略）}}]
			忽视 $a\ge b$ 的条件，认为 $a-b$ 可能为负而不敢使用。
		\item[\textbf{\textcolor{CTrap}{正确理解（条件前提）：}}]
			问题中已给定 $a\ge b$，故 $a-b\ge0$，可以直接用于表示小正方形边长。
		\item[\textbf{\textcolor{CTrap}{错因分析（前提遗忘）：}}]
			未注意题目中 $a\ge b$ 的条件，导致对表达式符号产生疑虑。
	\end{description}

\end{trapbox}
